\chapter{Vectors Tutorial}
\label{appendix:vectors}
\begin{quote}
  ``Are you making progress on these problems? Or are you making Congress?''
  \\ --Zuming Feng
\end{quote}

\section{Definitions}
Our definition of a vector is just an arrow that points from some reference point, denoted $\vec 0$, to another point.  In this paper, we are only considering vectors in the plane (e.g. two-dimensional vectors), but the concept extends easily to multiple dimensions.  It is assumed that the reader is familiar with vector addition, the concept of magnitude, etc.

In order to make the theory of vectors applicable to geometry, we can associate every point with the vector that points to it.

Of particular use in the world of vectors is the dot product defined by
\begin{definition}
  Let $\vec u$ and $\vec v$ be vectors.  Then the \textit{dot product} $\vec u \cdot \vec v$ is defined as
  \[ \vec u \cdot \vec v = |\vec u||\vec v| \cos \theta \]
  where $\theta$ is the angle between the two vectors.
\end{definition}

The dot product has the following very nice properties:
\begin{enumerate}
  \item $\vec u \cdot \vec u = |u|^2$.
  \item $\vec u \cdot \vec v = 0$ if and only if the vectors are perpendicular.
  \item The dot product is commutative, distributive, and associative.
\end{enumerate}

\begin{exercise}
  Let $A,B,C,D$ be points in the plane.  Prove that $AC \perp BD$ if and only if $(\vec A - \vec C) \cdot (\vec B - \vec D) = 0$.
\end{exercise}

\section{Triangle ABC}
Let $ABC$ be a triangle in the plane with circumcenter $O$.  Let $\vec O = \vec 0$ be the reference point.  Then it's obvious that $|\vec A| = |\vec B| = |\vec C| = R$, where $R$ is the circumradius of the triangle.

Now for the key fact:
\begin{fact}
  $\vec A \cdot \vec B = R^2 - \frac{c^2}{2}$.  Cyclic variations hold.
\end{fact}
\begin{proof}
  Using the definition,
  \begin{align*}
    \vec A \cdot \vec B &= | \vec A || \vec B | \cos \angle AOB \\
    &= R^2 \cos 2\angle ACB \\
    &= R^2 (1 - 2\sin^2 \angle ACB) \\
    &= R^2 (1 - 2\sin^2 C) \\
    &= R^2 - \half (2R \cdot \sin C)^2 \\
    &= R^2 - \frac{c^2}{2}
  \end{align*}
  the last step following by the Law of Sines.
\end{proof}

Barycentric coordinates allow us to write any point in the plane in the form $x \vec A + y \vec B + z \vec C$.  In that case, displacement vectors, points, etc. can all be written as a linear combination of $\vec A$, $\vec B$ and $\vec C$.  The above fact allows us to compute arbitrary dot products via the distributive law.  Using properties 1 and 2 above, we can construct criterion for perpendiculars and lengths; this leads to EFFT and the distance formula.

\section{Special Points}
Again setting $\vec O = \vec 0$, we have the orthocenter\footnote{To check that this is correct, verify that $(\vec H - \vec A) \cdot (\vec B - \vec C) = 0$, implying $HA \perp BC$.} $\vec H = \vec A + \vec B + \vec C$, the centroid $\vec G = \frac{1}{3} \vec H$, and the nine-point center $\vec N = \half \vec H$.  In other words, the points on the Euler Line are very, very nice.

This allows us, for example, to compute $OH^2$ as
\begin{align*}
  OH^2 &= \vec H \cdot \vec H \\
  &= (\vec A + \vec B + \vec C) \cdot (\vec A + \vec B + \vec C) \\
  &= 3R^2 + 2(R^2-a^2/2) + 2(R^2-b^2/2) + 2(R^2-c^2/2) \\
  &= 9R^2-a^2-b^2-c^2
\end{align*}


\chapter{Formula Sheet}
\label{appendix:formulasheet}
\begin{quote}
  My formula sheets were pretty thorough, but perhaps they were missing something.
  \\ --Richard Rusczyk
\end{quote}

%import tech
\section{Standard Formulas}

\makeatletter
\let\actuallabel\label
\makeatother
\renewcommand{\label}[1]{
%Do Nothing
}
\InputTempFile{tech}
\renewcommand{\label}[1]{\label{#1}}


\section{More Obscure Formulas}
Here's some miscellaneous formulas and the like.  These were not included in the main text.

From \cite{leibnizthm},
\begin{theorem}[Leibniz Theorem]
  Let $Q$ be a point with homogeneous barycentric coordinates $(u: v: w)$ with respect to $\triangle ABC$.  For any point $P$ on the plane $ABC$ the following relation holds:
  \[ uPA^2 + vPB^2 + wPC^2 = (u + v + w)PQ^2 + uQA^2 + vQB^2 + wQC^2 \]
\end{theorem}

\subsection{Other Special Points}
\begin{tabular}{r|l}
  Point & Coordinates \\ \hline
  Gregonne Point \cite{wolfram} & $\text{Ge} = ((s-b)(s-c) : (s-c)(s-a) : (s-a)(s-b))$ \\
  Nagel Point \cite{wolfram} & $\text{Na} = (s-a : s-b : s-c)$ \\
  Isogonal Conjugate \cite{infinity} & $P^\ast = \left( \frac{a^2}{x} : \frac{b^2}{y} : \frac{c^2}{z} \right)$ \\
  Isotomic Conjugate \cite{infinity} & $P^t = \left( \frac{1}{x} : \frac{1}{y} : \frac{1}{z} \right)$ \\
  Feuerbach Point \cite{fpoint} & $F = \left((b+c-a)(b-c)^2 : (c+a-b)(c-a)^2 : (a+b-c)(a-b)^2 \right)$ \\
  Nine-point Center & $N = (a\cos(B-C) : b\cos(C-A) : c\cos(A-B))$
\end{tabular}

\subsection{Special Lines and Circles}

\begin{tabular}{rl}
  Nine-point Circle & $-a^2yz-b^2xz-c^2xy+ \half (x+y+z)(S_Ax+S_By+S_Cz) = 0$ \\
  Incircle & $-a^2yz - b^2zx  - c^2xy + (x+y+z)\left( (s-a)^2x + (s-b)^2y + (s-c)^2z \right) = 0$ \\
  $A$-excircle \cite{book} & $-a^2yz - b^2zx - c^2xy + (x+y+z)\left( s^2 x + (s-c)^2 y + (s-b)^2 z \right) = 0$ \\
  Euler Line \cite{book} & $S_A(S_B-S_C)x + S_B(S_C-S_A)y + S_C(S_A-S_B)z = 0$
\end{tabular}

